% =====================================================================
%  CHAPTER 2: 食物中的生物活性成分
% =====================================================================
\chapter{食物中的生物活性成分}
\setcounter{chapter}{2}
\chapterintro{
\keyword{掌握}植物化学物的概念、分类和生物学作用；掌握多酚类化合物、有机硫化物、植物雌激素的生物学作用；掌握SPL（特定建议值）的概念及其与营养素DRIs的区别。\keyword{熟悉}类胡萝卜素、植物固醇的作用；生物活性物质与相关疾病。\keyword{了解}其他生物活性成分。
}

% ----- 第一节 植物化学物概述 -----
\section{植物化学物概述}

\begin{keybox}
\textbf{植物化学物（Phytochemicals）：}来自植物性食物的生物活性成分，是植物能量代谢过程中产生的多种中间或末端低分子量次级代谢产物。除个别是维生素的前体物（$\beta$-胡萝卜素）外，其余均为非传统营养素成分。

\textbf{分类（10类）：}①类胡萝卜素；②植物固醇；③皂苷类化合物；④芥子油苷；⑤多酚类化合物；⑥蛋白酶抑制剂；⑦单萜类；⑧植物雌激素；⑨有机硫化物；⑩植酸。

\textbf{生物活性（6种）：}①抑制肿瘤；②抗氧化；③免疫调节；④抑制微生物；⑤降低胆固醇；⑥其他：调节血压血糖、血小板血凝，抑制炎症等。
\end{keybox}

\begin{defbox}
\textbf{生物活性成分（Bioactive Food Components）：}食物中除了营养素以外的，对人体健康有益的微量非营养素成分。过去曾被称为"非营养素活性成分"，现已统一称为"食物中的生物活性成分"。按来源可分为植物来源（植物化学物）和动物来源两大类。
\end{defbox}

% ----- 第二节 各类植物化学物分述 -----
\section{各类植物化学物分述}

\subsection{类胡萝卜素（Carotenoids）}

\begin{defbox}
\textbf{类胡萝卜素：}广泛存在于果蔬中的脂溶性植物色素，迄今已发现700余种，约50种存在于人类膳食中。分为两大类：
\begin{itemize}[leftmargin=*]
    \item \textbf{胡萝卜素类（Carotenes）：}纯碳氢化合物，不含氧。包括$\alpha$-胡萝卜素、$\beta$-胡萝卜素、番茄红素（Lycopene）。
    \item \textbf{叶黄素类（Xanthophylls）：}含氧衍生物。包括叶黄素（Lutein）、玉米黄素（Zeaxanthin）、$\beta$-隐黄质。
\end{itemize}

\textbf{主要作用：}抗氧化（番茄红素淬灭单线态氧能力最强，是$\beta$-胡萝卜素的2倍以上）；维生素A原活性（$\beta$-胡萝卜素最高，1分子$\rightarrow$2分子视黄醇；番茄红素、叶黄素、玉米黄素无此活性）；抗肿瘤；免疫调节；视觉保护（叶黄素和玉米黄素构成视网膜黄斑色素，滤过蓝光，降低年龄相关性黄斑变性AMD风险）。
\end{defbox}

\subsection{植物固醇（Phytosterols）}

\begin{defbox}
\textbf{植物固醇：}存在于植物细胞膜中的固醇类化合物，结构与胆固醇类似（均为环戊烷多氢菲骨架，C17侧链有差异）。人体不能合成，全部来自膳食。主要种类：\keyword{$\beta$-谷固醇}（含量最丰富，约占50\%\~{}65\%）、豆固醇、菜油固醇。

\textbf{降胆固醇作用（最明确的核心功能）：}与胆固醇在肠道竞争进入混合微胶粒，减少胆固醇的溶解和吸收。每日摄入2\~{}3g可使血清LDL降低5\%\~{}15\%，是公认的辅助降脂功能性成分。\textbf{食物来源：}植物油、坚果和种子、豆类、全谷物。
\end{defbox}

\subsection{皂苷类化合物（Saponins）}

皂苷是一类具有表面活性（在水中振摇可产生持久泡沫，类似肥皂）的糖苷类化合物，由苷元（甾体或三萜类）和糖组成。

\textbf{主要来源：}大豆、鹰嘴豆、花生、菠菜、芦笋、燕麦、人参、甘草。

\textbf{主要作用：}调节脂质代谢（与胆固醇和胆汁酸结合降低肠道吸收）；抗菌/抗病毒（破坏微生物细胞膜）；抗肿瘤；抗血栓形成；免疫调节（人参皂苷可激活巨噬细胞和NK细胞）。注意：皂苷具有溶血活性，但经口摄入后胃肠道吸收极少，通常不造成全身性危害。

\subsection{芥子油苷（Glucosinolates）}

芥子油苷（硫代葡萄糖苷）主要存在于\keyword{十字花科蔬菜}（西兰花、卷心菜、白菜、甘蓝、萝卜、芥菜等）。植物细胞破碎后经黑芥子酶（Myrosinase）水解，生成具有生物活性的\keyword{异硫氰酸盐（ITCs）}，如\keyword{萝卜硫素（Sulforaphane）}。

\textbf{抗肿瘤作用——核心功能：}萝卜硫素是已知最强的天然Nrf2通路激活剂之一，强力诱导II相解毒酶（GST、UGT、NQO1）表达；流行病学证据支持降低多种癌症风险（肺癌、结肠癌、前列腺癌、乳腺癌）。此外还有调节炎症反应（抑制NF-$\kappa$B通路）和抗菌（抑制幽门螺杆菌）作用。

\subsection{多酚类化合物（Polyphenols）}

\begin{keybox}
\textbf{多酚类化合物：}植物性食物中含量最丰富、研究最深入的植物化学物类别。基本结构为2-苯基色原酮核心骨架，即一个或多个芳香环连接多个酚羟基。按碳骨架结构分为七大类：

\begin{enumerate}[leftmargin=*]
    \item \textbf{黄酮及黄酮醇类：}代表为\keyword{槲皮素}（膳食中最丰富的类黄酮之一）、芦丁、山奈酚。来源：洋葱、西兰花、蓝莓、苹果、红茶。
    \item \textbf{黄烷酮类：}代表为柚皮素、橙皮素。来源：柑橘类水果。
    \item \textbf{黄烷醇类（儿茶素类）：}代表为\keyword{儿茶素}、表儿茶素、EGCG。来源：绿茶（含量最丰富）、巧克力、苹果。
    \item \textbf{异黄酮类：}代表为大豆苷、\keyword{染料木素}、葛根素。来源：大豆及其制品。
    \item \textbf{双黄酮类：}代表为银杏黄酮。来源：银杏叶。
    \item \textbf{花色素类：}代表为矢车菊素、天竺葵素。来源：浆果、紫薯、紫甘蓝。
    \item \textbf{查尔酮类：}来源：苹果、甘草。
\end{enumerate}

\textbf{主要生物学作用：}①抗氧化——膳食最重要的抗氧化活性成分来源，直接清除自由基并通过激活Nrf2/ARE通路诱导抗氧化酶（SOD、CAT、GPx、HO-1）；②抗肿瘤——茶多酚（EGCG）和大豆异黄酮证据最充分（诱导II相解毒酶、抑制I相代谢酶、阻滞细胞周期、诱导凋亡、抑制血管生成）；③心血管保护——改善内皮功能、降低血压、抑制LDL氧化、抗血小板聚集；④抗炎——抑制NF-$\kappa$B通路和COX-2、iNOS；⑤调节肠道菌群——促进有益菌增殖，抑制致病菌。
\end{keybox}

\subsection{蛋白酶抑制剂（Protease Inhibitors）}

广泛存在于植物（特别是豆类和谷物）中的一类能与蛋白水解酶结合并抑制其活性的蛋白质或多肽。典型代表：大豆胰蛋白酶抑制剂（Bowman-Birk抑制剂BBI；Kunitz胰蛋白酶抑制剂KTI）。主要作用：抗肿瘤（抑制肿瘤相关蛋白酶活性、抑制自由基生成、调节癌基因表达）；免疫调节（BBI可抑制炎症性蛋白酶如弹性蛋白酶）。

蛋白酶抑制剂历来被视为"抗营养因子"（抑制胰蛋白酶、降低蛋白质消化率），但充分加热可使其彻底灭活，且现代研究发现在低浓度下具有抗癌防炎的健康效益。

\subsection{单萜类（Monoterpenes）}

由两个异戊二烯单元（C$_{10}$骨架）组成的萜类化合物，是植物精油的主要成分。代表性物质：紫苏醇（Perillyl Alcohol）、柠檬烯（Limonene）。主要来源：柑橘类水果果皮（柠檬烯）、薄荷油、香菜籽。

主要作用：抗肿瘤——紫苏醇和柠檬烯可抑制Ras蛋白异戊二烯化（Ras为GTP结合癌蛋白，需法尼基化才能定位到细胞膜发挥致癌效应），在动物模型中显示抑制乳腺、肝脏、胰腺肿瘤。

\subsection{植物雌激素（Phytoestrogens）}

\begin{keybox}
\textbf{植物雌激素：}植物来源的、能与哺乳动物雌激素受体（ER）结合并发挥类雌激素或抗雌激素效应的非甾体化合物。核心特征为\imp{双向调节作用}——低雌激素水平时激动ER（雌激素样作用），高雌激素水平时拮抗ER（竞争性抑制内源性雌二醇），使其成为SERM（选择性雌激素受体调节剂）的天然类似物。

\textbf{分类与来源：}
\begin{enumerate}[leftmargin=*]
    \item \textbf{异黄酮类（Isoflavones）：}大豆苷、染料木素、大豆黄素。\imp{主要来源：大豆及其制品}（豆粉、豆腐、豆浆），大豆中异黄酮含量约为0.1\%\~{}0.5\%。
    \item \textbf{木酚素类（Lignans）：}开环异落叶松脂素、罗汉松脂素。在肠道菌群作用下转化为肠内酯和肠二醇活性形式。\imp{亚麻籽是木酚素含量最高的食物}，也含于全谷物、种子类。
    \item \textbf{香豆素类（Coumestans）：}香豆雌酚。来源：发芽豆类（苜蓿芽、绿豆芽）。
    \item \textbf{芪类（Stilbenes）：}代表为\keyword{白藜芦醇}（Resveratrol）。来源：葡萄/葡萄酒、花生、虎杖。
\end{enumerate}

\textbf{主要生物学作用：}①预防骨质疏松——通过SERM-like作用促进成骨细胞活性、抑制破骨细胞分化；②改善围绝经期症状（缓解潮热、盗汗等）；③抗氧化；④心血管保护；⑤抗肿瘤——高大豆摄入人群（亚洲）的乳腺癌、前列腺癌发病率较低。
\end{keybox}

\subsection{有机硫化物（Organosulfur Compounds）}

\begin{keybox}
\textbf{有机硫化物：}主要存在于\keyword{葱属植物（Allium）}中的含硫植物化学物，以大蒜含量最为丰富，是其特征风味物质的来源。典型成分：蒜素（Allicin）——完整蒜瓣切碎或压碎后由蒜氨酸（Alliin）经蒜氨酸酶催化生成，在室温下不稳定，可进一步分解为二烯丙基二硫化物（DADS）、二烯丙基三硫化物（DATS）、阿霍烯（Ajoene）等。

\textbf{主要生物学作用：}
\begin{enumerate}[leftmargin=*]
    \item \textbf{抗菌作用：}蒜素具有广谱抗菌活性，杀菌效力约相当于1\%青霉素，对革兰阳性和阴性细菌、真菌、寄生虫均有抑制作用。细菌不易对蒜素产生耐药性。
    \item \textbf{抗氧化作用：}直接清除自由基，并诱导抗氧化酶的表达。
    \item \textbf{调节血脂：}降低血清TC、TG和LDL，升高HDL。
    \item \textbf{抗血栓形成：}抑制血小板聚集，有一定纤溶促进作用（类似阿司匹林但较弱）。
    \item \textbf{抗肿瘤作用：}阻断\imp{亚硝胺}的内源性合成——这是大蒜防癌的重要机制之一；诱导II相解毒酶表达，加速致癌物灭活和排泄。
\end{enumerate}

\textbf{食用建议：}蒜素在高温下易被破坏，建议大蒜切碎后室温放置10\~{}15分钟（让蒜氨酸酶充分作用生成蒜素），再用于凉拌或低温快炒。
\end{keybox}

\subsection{植酸（Phytic Acid, IP6）}

\begin{defbox}
\textbf{植酸（肌醇六磷酸, IP6）：}豆类、全谷物和坚果中天然存在的肌醇磷酸酯。含有六个磷酸基团，对二价矿物离子（Ca$^{2+}$、Fe$^{2+}$/Fe$^{3+}$、Zn$^{2+}$、Mg$^{2+}$）具有强烈的螯合能力。

\textbf{认识的历史演变：}
\begin{enumerate}[leftmargin=*]
    \item \textbf{传统认识（抗营养因子）：}在肠道螯合矿物质，降低矿物质的生物利用度，是发展中国家以谷物为主食人群矿物质缺乏（尤其是铁、锌）的重要影响因素。
    \item \textbf{现代认识（有益生物活性成分）：}具有抗氧化（铁螯合阻断Fenton反应产生自由基）和抗肿瘤（抑制增殖、诱导分化）的生物学作用，被认为是潜在的抗癌活性物质。
\end{enumerate}
\end{defbox}

% ----- 第三节 生物活性物质与相关疾病 -----
\section{生物活性物质与相关疾病}

\begin{defbox}
\textbf{植物化学物与慢性病防治：}大量流行病学研究和临床试验表明，富含植物化学物的膳食模式（多蔬果、全谷物、大豆、坚果）与多种慢性非传染性疾病（NCD）的风险降低显著相关。

\textbf{主要作用机制：}
\begin{itemize}[leftmargin=*]
    \item \textbf{抗肿瘤：}多途径协同——诱导II相解毒酶（萝卜硫素、蒜素）、抑制I相致癌物代谢酶（多酚类）、阻滞细胞周期（多酚类）、诱导凋亡（多酚类、皂苷）、抑制血管生成（多酚类）、阻断致癌物合成（有机硫化物阻断亚硝胺合成）、抑制Ras蛋白异戊二烯化（单萜类）。
    \item \textbf{心血管保护：}降胆固醇（植物固醇竞争吸收、皂苷结合胆汁酸、有机硫化物抑制合成酶）、抗氧化/抗炎（多酚类抑制LDL氧化和NF-$\kappa$B通路）、抗血小板聚集（有机硫化物、皂苷）、改善内皮功能。
    \item \textbf{代谢调节：}调节血糖（多酚类、硫辛酸）、改善胰岛素敏感性。
    \item \textbf{骨骼健康：}植物雌激素通过SERM样作用抑制骨丢失。
    \item \textbf{视觉保护：}叶黄素/玉米黄素保护视网膜，降低AMD风险。
\end{itemize}
\end{defbox}

% ----- 第四节 SPL -----
\section{特定建议值（SPL）}

\begin{keybox}
\textbf{特定建议值（Specific Proposed Levels, SPL）：}SPL是DRIs指标体系的第七项，专门针对\textbf{非营养素生物活性成分}（主要为植物化学物）而设立。其定义为：以降低成年人膳食相关非传染性慢性病（NCD）风险为目标的\textbf{其他膳食成分}每日摄入量。

\textbf{SPL与营养素DRIs的本质区别：}
\begin{enumerate}[leftmargin=*]
    \item \textbf{适用对象不同：}SPL针对非必需的非营养素活性成分（植物化学物），DRIs的EAR/RNI/AI/UL针对必需营养素。
    \item \textbf{底层逻辑不同：}SPL回答"摄入多少有助于降低慢病风险"（促健康导向），RNI回答"需要多少满足基本生理需求"（防缺乏导向），UL回答"多少以内安全"。
    \item \textbf{不存在EAR和RNI：}植物化学物为非必需、非缺乏性物质，因此无EAR（无"满足50\%个体"的概念）和RNI。
\end{enumerate}

\textbf{已提出SPL的植物化学物举例（中国DRIs 2023版）：}
\begin{itemize}[leftmargin=*]
    \item \textbf{大豆异黄酮：}SPL以苷元计，基于改善围绝经期症状和预防骨质疏松的证据。
    \item \textbf{植物固醇：}SPL基于降低血清LDL胆固醇的证据。
    \item \textbf{原花青素、花色苷、番茄红素、叶黄素等：}亦有相应SPL值被提出或正在论证中。
\end{itemize}

\textbf{意义：}SPL的设立标志着营养科学从"预防缺乏"向"促进健康和预防慢病"的范式转变，承认了植物化学物在膳食营养指导中的独立地位，为膳食指南中"多吃蔬果、全谷物、大豆"的建议提供了剂量化依据。
\end{keybox}

% ----- 第五节 动物来源的生物活性成分 -----
\section{动物来源的生物活性成分（了解）}

\begin{itemize}[leftmargin=*]
    \item \textbf{辅酶Q（泛醌，CoQ10）：}呼吸链电子传递体，参与ATP合成；脂溶性抗氧化剂（保护细胞膜和LDL）；心血管保护。
    \item \textbf{硫辛酸：}"万能抗氧化剂"兼有水溶性和脂溶性，能透过血脑屏障。再生GSH和维生素C/E；改善胰岛素敏感性（II型糖尿病辅助治疗）；神经保护。
    \item \textbf{褪黑素：}主要由松果体合成。调节生物节律（昼夜节律和睡眠周期）；强效抗氧化剂（穿越所有生物屏障）。
    \item \textbf{$\gamma$-氨基丁酸（GABA）：}抑制性神经递质，降血压、镇静安神。
    \item \textbf{左旋肉碱（L-Carnitine）：}脂肪酸$\beta$-氧化的转运载体（将长链脂肪酸转运至线粒体基质）。来源：红肉、奶制品。
\end{itemize}

% ----- 复习题 -----
\reviewcontent{
\section{复习题}

\begin{enumerate}[leftmargin=*, itemsep=8pt]
    \item \textbf{【名词解释】}植物化学物（Phytochemicals）；植物雌激素的双向调节作用；SPL（特定建议值）
    \item \textbf{【简答题】}列出植物化学物的10个分类。
    \item \textbf{【简答题】}列出植物化学物的6种生物活性。
    \item \textbf{【简答题】}简述多酚类化合物的七大类别及其主要生物学作用。
    \item \textbf{【简答题】}简述有机硫化物的主要生物学作用及食用建议。
    \item \textbf{【简答题】}试述植物雌激素的分类、来源及其双向调节机制。
    \item \textbf{【简答题】}简述SPL与营养素DRIs（RNI、UL）的本质区别。
    \item \textbf{【简答题】}简述植酸从"抗营养因子"到"潜在抗癌活性物质"的认识转变。
    \item \textbf{【论述题】}试述植物化学物在慢性非传染性疾病（NCD）防治中的作用及机制。
\end{enumerate}

\subsection*{参考答案要点}

\begin{enumerate}[leftmargin=*, itemsep=5pt]
    \item \textbf{植物化学物：}来自植物性食物的生物活性成分，是植物能量代谢过程中产生的多种中间或末端低分子量次级代谢产物。除个别是维生素的前体物（$\beta$-胡萝卜素）外，其余均为非传统营养素成分。\textbf{双向调节作用：}低雌激素水平时激动ER（雌激素样作用），高雌激素水平时拮抗ER（竞争性抑制内源性雌二醇）。\textbf{SPL：}DRIs第七项指标，针对非营养素生物活性成分，以降低成人膳食相关NCD风险为目标的每日摄入量。
    \item ①类胡萝卜素；②植物固醇；③皂苷类化合物；④芥子油苷；⑤多酚类化合物；⑥蛋白酶抑制剂；⑦单萜类；⑧植物雌激素；⑨有机硫化物；⑩植酸。
    \item ①抑制肿瘤；②抗氧化；③免疫调节；④抑制微生物；⑤降低胆固醇；⑥其他：调节血压血糖、血小板血凝，抑制炎症等。
    \item \textbf{七大类：}①黄酮及黄酮醇类（槲皮素）；②黄烷酮类（柚皮素）；③黄烷醇类/儿茶素类（EGCG）；④异黄酮类（染料木素）；⑤双黄酮类（银杏黄酮）；⑥花色素类（矢车菊素）；⑦查尔酮类。\textbf{主要作用：}抗氧化（激活Nrf2/ARE通路诱导抗氧化酶）、抗肿瘤（诱导II相解毒酶、诱导凋亡、抑制血管生成）、心血管保护、抗炎（抑制NF-$\kappa$B）、调节肠道菌群。
    \item \textbf{作用：}①抗菌——蒜素广谱抗菌，效力约相当于1\%青霉素，不易产生耐药性；②抗氧化；③调脂——降低TC/TG/LDL，升高HDL；④抗血栓——抑制血小板聚集；⑤抗肿瘤——阻断亚硝胺内源性合成、诱导II相解毒酶。\textbf{食用建议：}大蒜切碎后室温放置10\~{}15分钟（让蒜氨酸酶充分作用生成蒜素），避免高温长时间烹饪。
    \item \textbf{分类：}异黄酮类（大豆苷、染料木素）、木酚素类（亚麻籽含量最高，经肠道菌群转化为肠内酯/肠二醇）、香豆素类（香豆雌酚，发芽豆类）、芪类（白藜芦醇，葡萄/葡萄酒）。\textbf{机制：}低雌激素水平时结合ER发挥雌激素样作用，高雌激素水平时竞争性抑制内源性雌二醇发挥抗雌激素效应，是天然的SERM。作用包括预防骨质疏松、改善围绝经期症状、抗肿瘤。
    \item ①SPL针对非必需的非营养素（植物化学物），RNI/UL针对必需营养素；②SPL回答"摄入多少有助降低慢病风险"，RNI回答"需要多少满足生理需求"，UL回答"多少以内安全"；③SPL无EAR（无需满足50\%个体），仅基于流行病学剂量-反应关系；④SPL标志着从"预防缺乏"到"促健康防慢病"的范式转变。
    \item \textbf{传统认识：}螯合二价矿物质（Ca$^{2+}$、Fe$^{2+}$/Fe$^{3+}$、Zn$^{2+}$、Mg$^{2+}$），降低矿物质生物利用度，是发展中国家以谷物为主食人群矿物质缺乏的重要影响因素（抗营养因子）。\textbf{现代认识：}具有抗氧化（铁螯合阻断Fenton反应产生自由基）和抗肿瘤（抑制增殖、诱导分化）的生物学作用，被认为是潜在的抗癌活性物质（有益生物活性成分）。
    \item \textbf{抗肿瘤：}多途径协同——诱导II相解毒酶（萝卜硫素激活Nrf2、有机硫化物）、抑制I相致癌物代谢酶（多酚类）、阻滞细胞周期和诱导凋亡（多酚类、皂苷类）、抑制血管生成（多酚类）、阻断致癌物合成（有机硫化物阻断亚硝胺）、抑制Ras蛋白异戊二烯化（单萜类）。\textbf{心血管保护：}降胆固醇（植物固醇竞争吸收、皂苷结合胆汁酸、有机硫化物抑制合成酶）、抗氧化抗炎（多酚类抑制LDL氧化和NF-$\kappa$B）、抗血小板聚集（有机硫化物、皂苷）、改善内皮功能。\textbf{其他：}植物雌激素预防骨质疏松（SERM样作用）、叶黄素/玉米黄素保护视觉、多酚类调节肠道菌群等。
\end{enumerate}

\newpage
}
